\section{Discussion}

One supported result is about model structure: aligning the estimand and
capacity with the decision improved a learned controller relative to a dense
residual.  CFCMT predicts a normalized ordering among phases from the same
state, keeps a rigid local anchor and allows an all-pair correction only at a
weight selected on complete held-out target seeds.  This reduced external
offline regret and closed-loop waiting relative to a same-information dense
residual.  Because the dense comparison also receives matched action
contrasts, the gain cannot be explained by reference subtraction alone.

A distinct claim about transferred data is conditional rather than uniformly
positive or negative.  The v91 comparator shares the target labels, feature
semantics, model family and executor with CFCMT, and assigns source-labelled
samples zero weight in deployed scoring.  However, its retained source prior
can resolve exact prediction ties, so it is a legacy near-target-only
comparator rather than a strict source-free arm.  Across 64 fresh closed-loop
seeds, CFCMT reduced Jinan waiting by 3.045\%, with all 64 seed-level city
contrasts favourable, but increased Los Angeles waiting by 7.865\%, with only
22 of 64 contrasts favourable.  The equal-city mean therefore rejected a
general cross-city source-contribution claim and exposed negative transfer as
a network-dependent failure mode.

A separate post-v91 target-offline selector then made the network dependence
operational before its fresh rollouts: it returned exact target-only control
for Los Angeles and admitted a Hangzhou source component at mass one for
Jinan.  On 56 further Jinan seeds (336 rollouts), the admitted controller
reduced waiting by 4.527\% relative to a strict zero-source-row target model
(paired 95\% interval, 4.279--4.787\%); all 56 contrasts were favourable
($p=7.5475\times10^{-11}$, Wilcoxon signed-rank test).  Because the abstention
and admission decisions preceded these rollouts, this is evidence for a
preselected Jinan controller-pair gain rather than a post-hoc deletion of Los
Angeles.
The comparison still does not isolate source rows at fixed architecture: the
strict target comparator has lower capacity, and both arms use the same risk
multiplier of 0.5 inherited from earlier Jinan closed-loop development.  It
also does not establish superiority over phase pressure.

The experiments also explain why the earlier full mechanism-world-model route
was not retained as the final controller.  MC-WM decomposed queue propagation,
red accumulation, mobility, spillback and control cost.  That decomposition is
useful for diagnosis, but each auxiliary prediction introduces error before the
model answers the narrower control question.  Good next-state or mechanism loss
does not guarantee the correct phase ordering, and flexible latent or graph
paths can dominate the rigid signal under limited target support.  The final
method is therefore a reverse design from the successful rigid model: preserve
declared parents and local mechanism features, then add only an antisymmetric
ranking correction whose effect is measurable and bounded.  This is a change
in scientific hypothesis, not a residual patch to rescue a predetermined
``full'' architecture.

Pressure control remains the operational reference.  MaxPressure and phase
pressure were better than CFCMT in both external cities, with lower system
vehicle-hours and slightly higher completion.  Traffic-signal control has a
strong local-rule ceiling because queue pressure directly reflects service
imbalance.  A learned method that beats fixed time or a dense neural baseline
but not pressure control is not a replacement controller.  CFCMT is better
positioned as a transfer layer for settings where the decision depends on
effects absent from a pressure score, or as a structured model used alongside
rather than rhetorically above a rule.  The discarded pressure-guard
experiments in the supplement reinforce this point, but pressure is not hidden
inside the final v43 policy.

The absence of same-protocol CrossLight, X-Light, MetaLight and GESA
reproductions is also a substantive comparison gap, not merely a citation
issue.  Their published protocols use different target interaction budgets,
state and action interfaces, simulators and training objectives.  Reporting
their native numbers beside the paired SUMO matrix would manufacture a ranking
rather than estimate one.  A broader benchmark should reimplement each method
behind the same safe phase executor and expose three separate target budgets:
zero labels, passive target transitions and matched counterfactual target
actions.

Target information changes the interpretation.  The target simulator is
uncalibrated against field trajectories, which reduces calibration effort, but
its action outcomes are labelled target data.  The external result is thus a
few-shot/full-budget offline-adaptation result with 165 and 911 groups, not a
demonstration that a source-city policy can be deployed without target
transitions.  A genuine zero-shot arm remains scientifically useful as the
left endpoint of an information curve, but it is not the basis of the present
closed-loop claim.  Conversely, field loop-detector or trajectory data could
adapt passive state representations without supplying the counterfactual
outcome of an unexecuted phase.

The source-comparison results also change how the adaptation budget should be
used.
More source data is not monotonically beneficial after target labels are
available.  Initial target-seed-blocked gates either retained city regressions
or removed them by selecting target only in all seven development cities
(Supplementary Table~\ref{tab:source-gate-development}).  A subsequent remedy
ladder tested mechanism-specific priors, inverse-budget shrinkage, local
right-of-way context, state-conditioned downstream utility, leave-city-out
meta-utility and source-only or target-calibrated hierarchical priors
(Supplementary Table~\ref{tab:source-remedy-ladder}).  Forced source arms still
showed useful headroom, including improvement in five of seven city means for
the final target-calibrated prior, but its city interval crossed zero and no
source survived the complete source-identity controls.  The only state-level
utility gate that passed offline development produced a large Cologne
closed-loop regression before matrix expansion.  These experiments validate
exact abstention as a fallback and expose source-benefit identification as the
unresolved multicity problem; they do not extend the separate Jinan-specific
v98 result into a general deployable selector.

Several limitations constrain generalization.  First, the formal external
unit count is two cities.  Multiple Jinan demand files improve within-city
coverage but do not create independent cities, so all intervals are explicitly
descriptive.  Second, the Los Angeles and Jinan scenarios are CityFlow-derived
networks converted and semantically audited for SUMO; they are not calibrated
city-wide digital twins.  Third, training counterfactuals share the same SUMO
physics family as evaluation.  Cross-network transfer is real at the topology,
demand and phase-program levels, but sim-to-real transfer is not tested.
Fourth, the seven-policy closed-loop matrix uses two seeds, whereas v91 uses 64
seeds over the same two target geometries and v98 uses 56 further seeds on
Jinan only.  These protocols are reported separately rather than pooled.  All
use a fixed 3,600-s horizon.  Unfinished vehicles are included, yet longer
incidents, non-recurrent demand and detector failure remain outside the
experiment.

The sequential development record is another limitation and a strength.  A
diagnostic external seed exposed a one-city regression and motivated removal
of an arbitrary adaptation cap and stricter complete-seed cross-fitting.  That
diagnostic seed was then permanently excluded, and a new seed authorized the
final test.  This is not equivalent to a single untouched preregistration;
reporting the failed stages is necessary to expose researcher degrees of
freedom.  The final v43 cache, selectors, deployment models and closed-loop
matrix are nevertheless hash-linked and frozen before their respective next
stage.

A later execution-layer programme provides a stronger check on this boundary.
Its favourable 54-seed adaptive estimate shrank almost to zero on 64 newly
generated Jinan seeds and failed the preregistered bootstrap, sign-test and
worst-seed criteria.  Because that programme used a different 120-s latent
originator and global cooldown, it neither rescues nor invalidates the v43
anchored all-pair comparison.  It does show why adaptive search and independent
confirmation must not be pooled, and why the present evidence cannot support
learned-policy superiority over phase pressure.

The fresh v91 near-target-only test more directly rejects a uniform cross-city
source-benefit claim.  Unlike the execution-layer programme, it holds the
CFCMT model family, target adaptation budget and executor fixed.  Its retained
source-prior tie-break prevents a strict source-row causal interpretation, but
the opposite city effects cannot be dismissed as a different policy class.
Offline regret and short-horizon action ranking did not reliably predict the
closed-loop value of source data, particularly in the Los Angeles corridor
where one seed produced a large congestion cascade.

The post-v91 evidence narrows but does not close that attribution gap.  V98
shows that a frozen target-offline rule can abstain for Los Angeles and retain a
repeatable Jinan-specific controller gain against its declared lower-capacity
zero-source-row target model.
The development-only v123/v157a analysis reported favourable model-path
differences at five target-label budgets, but the target and source paths used
different Jinan domain handling.  These historical contrasts therefore do not
isolate source-row value.  V157B corrected the comparator at B100 only: on the
reused 22-seed Jinan selector, source minus target was $-0.01357$ (paired 95\%
interval $[-0.01784,-0.00935]$), with 20/22 seeds improving.  The source arm
still cost $+0.02534$ relative to phase pressure.  V157C then ran the fixed
three-seed Jinan native one-action protocol.  Seeds 80314 and 88625 completed,
but seed 27178 remained invalid because its fixed 480-s checkpoint had no
eligible, scorable traffic light; the seed was neither dropped nor replaced.
Across the two valid seeds, the descriptive source-minus-target mean was
$+0.00236$ (one of two seeds improved), source minus matched placebo was
$-0.00232$, target minus PhasePressure was $+0.00125$, and source minus
PhasePressure was $+0.00362$, with source worse in both seeds.  Recorded
collision counts were zero.  PhasePressure is the reference heuristic, not a
learned arm.  Because the fixed roster was incomplete, these values do not
support a three-seed inference or confirm the reused-selector advantage.
Moreover, only two of five effective source overrides improved their native
PhasePressure continuation, and the predicted source--target ordering agreed
with the native outcome in only one of four different-action windows.  The
limitation therefore extends beyond the missing seed to action-benefit ranking
on native continuations.  V158 made one final frozen attempt to address that
ranking mismatch.  It collected 100 native-prefix action groups from ten fixed
Jinan development seeds and fitted native residual calibrators in five
whole-seed folds.  Source minus target was $-0.00113$ (paired 95\% interval
$[-0.00240,+0.00019]$; 6/10 seeds improved), source minus matched placebo was
$-0.00072$ ($[-0.00219,+0.00091]$; 7/10), and source minus PhasePressure was
$+0.00014$ ($[-0.00239,+0.00261]$; 5/10).  All three preregistered OOF
comparisons failed, so the reserve was withheld.  This result authorizes no
controller adoption, unseen-city claim or post-hoc retuning; it closes the
native-prefix ranking-remediation route under the frozen Jinan sparse
controller.  No single experiment therefore combines
fixed architecture, strict source removal, fresh multicity closed-loop
confirmation and pressure-policy superiority.

Any future study would require a new precommitted design and independent city
units rather than a relaxed threshold or another sweep on the V158 outcomes.
It would need to freeze the link between compatibility and downstream action
ranking before evaluation on additional external cities and a calibrated
digital twin or prospective signal intervention.  Such a study could
test whether the local pairwise advantage survives simulator
mismatch and whether source mechanisms add value in objectives where pressure
is incomplete, including multimodal priority, spillback risk, emissions or
network-level coordination.  Until then, the supported claim is narrower:
causal parent restrictions and antisymmetric action ranking improve a
simulation-assisted residual relative to dense modelling, while a
predeployment rule retained a Jinan-specific controller-pair gain that does not
isolate source-row contribution.  Classical pressure remains the stronger
controller on the reported targets.
